\documentclass[12pt,a4paper]{article}
\usepackage[left=15mm,right=15mm,top=14mm,bottom=19mm,footskip=9mm]{geometry}
\usepackage{iftex}
\ifPDFTeX
  \usepackage[T5]{fontenc}
  \usepackage[utf8]{inputenc}
  \usepackage{mathptmx}
\else
  \usepackage{fontspec}
  \IfFontExistsTF{Times New Roman}{\setmainfont{Times New Roman}}{\setmainfont{Tinos}}
\fi
\usepackage{amsmath,amssymb}
\usepackage{tikz}
\usetikzlibrary{arrows.meta}
\usepackage[most]{tcolorbox}
\usepackage{fancyhdr,lastpage}
\usepackage{tabularx,array}
\definecolor{sectionpurple}{RGB}{155,42,143}
\definecolor{notebackground}{RGB}{229,240,245}
\definecolor{noteborder}{RGB}{255,55,55}
\definecolor{graphred}{RGB}{188,55,61}
\definecolor{graphblue}{RGB}{45,95,200}
\setlength{\parindent}{0pt}
\setlength{\parskip}{2pt}
\setlength{\headheight}{15pt}
\setlength{\abovedisplayskip}{5pt}
\setlength{\belowdisplayskip}{5pt}
\setlength{\abovedisplayshortskip}{3pt}
\setlength{\belowdisplayshortskip}{3pt}
\pagestyle{fancy}
\fancyhf{}
\renewcommand{\headrulewidth}{0pt}
\renewcommand{\footrulewidth}{0.4pt}
\fancyfoot[L]{\fontsize{8.5}{10}\selectfont Tài liệu tổng hợp theo TDMXX13 (Tư Duy Mở)}
\fancyfoot[R]{\fontsize{8.5}{10}\selectfont Trang \thepage/\pageref{LastPage} -- Bài giảng 06}
\newtcolorbox{notebox}{colback=notebackground,colframe=noteborder,
  boxrule=0.8pt,arc=2pt,left=7pt,right=7pt,top=3pt,bottom=3pt,
  before skip=3pt,after skip=3pt}
\newcommand{\partheading}[1]{\par\vspace{5pt}\noindent\textbf{\textcolor{sectionpurple}{#1}}\par\vspace{1pt}}
\newcommand{\example}[2][TDM21]{\par\vspace{3pt}\noindent\textbf{\textcolor{sectionpurple}{Ví dụ #2:} \textcolor{noteborder}{[#1]}}\enspace}
\newcommand{\tightitem}{\setlength{\itemsep}{1pt}\setlength{\topsep}{1pt}\setlength{\parsep}{0pt}\setlength{\partopsep}{0pt}}
\newcommand{\loigiai}[1][7mm]{\par\noindent\textbf{\textcolor{sectionpurple}{Giải:}}\par\vspace{#1}}
\begin{document}
\fontsize{12}{14.6}\selectfont
\begin{center}
  \textbf{BÀI GIẢNG 06}\\[-1pt]
  \textbf{CƠ BẢN VỀ NGUYÊN HÀM}\\[-2pt]
  {\fontsize{10.5}{12}\selectfont\itshape Biên soạn theo chủ đề TDMXX13 -- Tư Duy Mở}
\end{center}
\vspace{-7pt}

\partheading{PHẦN I. ĐỊNH NGHĨA NGUYÊN HÀM}
Cho hàm số $f(x)$ xác định trên $K$ ($K$ là khoảng hoặc đoạn hoặc nửa khoảng của $\mathbb{R}$).

\textbf{Định nghĩa.} Hàm số $F(x)$ được gọi là \textbf{nguyên hàm} của hàm số $f(x)$ trên $K$ nếu $F'(x)=f(x)$ với mọi $x\in K$.

\begin{notebox}
\example{01} Hãy tìm một nguyên hàm của hàm số:
\begin{enumerate}\tightitem
  \item $f(x)=3x^2$ trên $\mathbb{R}$.
  \item $f(x)=\dfrac1x$ trên khoảng $(0;+\infty)$.
\end{enumerate}
\end{notebox}
\loigiai[18mm]

Ta có: $F'(x)=f(x)\iff [F(x)+C]'=f(x)$.

\textbf{Nhận xét.} Nếu $F(x)$ là một nguyên hàm của hàm số $f(x)$ trên $K$ thì $F(x)+C$, $C\in\mathbb R$, là họ tất cả các nguyên hàm của $f(x)$ trên $K$.

\textbf{Kí hiệu:} $\displaystyle \int f(x)\,dx=F(x)+C$.

Suy ra:
\[
\bigl[F(x)+C\bigr]'=\left(\int f(x)\,dx\right)'
\iff F'(x)=\left(\int f(x)\,dx\right)'=f(x).
\]
\textit{Đạo hàm của nguyên hàm của một hàm số bằng chính hàm số đó.}

Lấy vi phân hai vế, ta được:
\[
 d\bigl(F(x)+C\bigr)=d\left(\int f(x)\,dx\right)
 \iff F'(x)\,dx=d\left(\int f(x)\,dx\right)=f(x)\,dx.
\]
\textit{Vi phân của nguyên hàm bằng biểu thức dưới dấu nguyên hàm.}

\begin{notebox}
\example{02} Hãy tính:
\begin{enumerate}\tightitem
  \item $\displaystyle\int\cos x\,dx$.
  \item $\displaystyle\int4t^3\,dt$.
  \item $\displaystyle\int\frac{1}{\cos^2s}\,ds$.
\end{enumerate}
\end{notebox}
\loigiai[21mm]

\clearpage
\partheading{PHẦN II. TÍNH CHẤT CỦA NGUYÊN HÀM}
Các tính chất cơ bản:
\begin{enumerate}\tightitem
 \item $\displaystyle\int f'(x)\,dx=f(x)+C$.
 \item $\displaystyle\int kf(x)\,dx=k\int f(x)\,dx$ \quad ($k$ là hằng số khác $0$).
 \item $\displaystyle\int\bigl[f(x)\pm g(x)\bigr]dx
       =\int f(x)\,dx\pm\int g(x)\,dx$.
\end{enumerate}
\textbf{Nhận xét:}\par\vspace{5mm}

\begin{notebox}
\example{01} Hãy tính:
\begin{enumerate}\tightitem
  \item $\displaystyle I=\int\bigl(2x-e^x+3\cos x\bigr)\,dx$.
  \item $\displaystyle J=\int\left(3x^2+1-\frac1x\right)\,dx$.
\end{enumerate}
\end{notebox}
\loigiai[45mm]

\partheading{PHẦN III. SỰ TỒN TẠI NGUYÊN HÀM}
\textbf{Chấp nhận kết quả.} Mọi hàm số $f(x)$ liên tục trên $K$ đều có nguyên hàm trên $K$.

\begin{notebox}
\example{01} Xét các ví dụ sau:
\begin{enumerate}\tightitem
  \item Hàm số $f(x)=x^{\frac12}$ liên tục trên khoảng $(0;+\infty)$ nên có nguyên hàm trên khoảng $(0;+\infty)$ và
  \[
    \int x^{\frac12}\,dx=\frac23x^{\frac32}+C.
  \]
  \item Hàm số $f(x)=\dfrac1{\cos^2x}$ liên tục trên từng khoảng
  $\left(-\frac\pi2+k\pi;\frac\pi2+k\pi\right)$, $k\in\mathbb Z$,
  nên có nguyên hàm trên từng khoảng đó và
  \[
     \int\frac1{\cos^2x}\,dx=\tan x+C.
  \]
\end{enumerate}
\end{notebox}

\clearpage
\partheading{PHẦN IV. BẢNG NGUYÊN HÀM THƯỜNG GẶP}
\begin{notebox}
\renewcommand{\arraystretch}{1.65}
\begin{tabularx}{\linewidth}{@{}>{\raggedright\arraybackslash}X >{\raggedright\arraybackslash}X@{}}
$\displaystyle\int0\,dx=C$ & $\displaystyle\int dx=x+C=\int1\,dx$\\
$\displaystyle\int x^\alpha\,dx=\frac{x^{\alpha+1}}{\alpha+1}+C$ & $\displaystyle\int\frac1x\,dx=\ln|x|+C$\\
\multicolumn{2}{@{}l@{}}{\hspace{1em}\footnotesize (Ở công thức lũy thừa: $\alpha\ne-1$, xét trên khoảng thích hợp.)}\\
$\displaystyle\int e^x\,dx=e^x+C$ & $\displaystyle\int a^x\,dx=\frac{a^x}{\ln a}+C$\\
$\displaystyle\int\cos x\,dx=\sin x+C$ & $\displaystyle\int\sin x\,dx=-\cos x+C$\\
$\displaystyle\int\frac{dx}{\cos^2x}=\tan x+C$ & $\displaystyle\int\frac{dx}{\sin^2x}=-\cot x+C$\\
\end{tabularx}
\vspace{1pt}
\footnotesize Điều kiện của hàm mũ: $a>0$, $a\ne1$. \textit{Hiệu đính theo toán học: bản ảnh nguồn ghi $+\cot x+C$ ở dòng cuối, thiếu dấu trừ.}
\end{notebox}

\textbf{Lưu ý.} Khi yêu cầu tính $\displaystyle\int f(x)\,dx$, ta hiểu là tính trên từng khoảng xác định của $f(x)$.

\begin{notebox}
\example{01} Hãy tính:
\begin{enumerate}\tightitem
  \item $\displaystyle A=\int\left(x^3+\frac3{\sqrt x}\right)\,dx$.
  \item $\displaystyle B=\int 2^{x+3}\,dx$.
  \item $\displaystyle C=\int\frac{(x^2+1)^2}{x^3}\,dx$.
  \item $\displaystyle D=\int(e^x+1)(e^{-x}+1)\,dx$.
\end{enumerate}
\end{notebox}
\loigiai[65mm]

\vspace{1pt}
{\centering\textbf{\textcolor{sectionpurple}{--- Hết ---}}\par}
\end{document}